\documentclass[12pt]{article}
\usepackage[margin=1in]{geometry}
\usepackage{amsmath} % the standard package for serious mathematics
\usepackage{amssymb} % additional symbols

\title{Typesetting Mathematics}
\author{Your Name}
\date{\today}

\begin{document}
\maketitle

\section{Inline versus display}

Math inside a sentence goes between dollar signs: the area of a circle
is $A = \pi r^2$. To set something on its own centred line, use a
display:
\[
    E = mc^2.
\]

\section{Building blocks}

Superscripts use \verb|^| and subscripts use \verb|_|. Group anything
longer than one character in braces --- write \verb|x^{10}|, not
\verb|x^10|:
\[
    x^{10}, \qquad a_{i+1}, \qquad \frac{-b \pm \sqrt{b^2 - 4ac}}{2a}.
\]

Sums and integrals grow to fit their contents:
\[
    \sum_{i=1}^{n} i = \frac{n(n+1)}{2}, \qquad
    \int_0^1 x^2 \, dx = \frac{1}{3}.
\]

\section{Numbered and aligned equations}

The \texttt{equation} environment adds an automatic number:
\begin{equation}
    a^2 + b^2 = c^2.
\end{equation}

To line several rows up at the equals sign, use \texttt{align}. The
\verb|&| marks the alignment point and \verb|\\| ends each row:
\begin{align}
    (x+1)^2 &= (x+1)(x+1) \\
            &= x^2 + 2x + 1.
\end{align}

\end{document}
